\documentclass[a4paper,11pt]{article}

\usepackage[ngerman]{babel}
\usepackage[top=3cm,bottom=3cm,left=3cm,right=3cm]{geometry} %NICHT IN KEMIE2


%\usepackage{microtype} %schöneres typesetting  mit [stretch=10] für nur 1% anstelle von 2%
%\usepackage{lmodern} %Latin Modern FONT
%\usepackage{pifont} %schönere Font



\usepackage{amsmath}
\usepackage{nicefrac}
\usepackage{amssymb}
%\usepackage{mathtools}

\usepackage{titlesec} %Um Section, etc bearbeiten zu können
\usepackage{setspace}
\usepackage{enumitem} %enumerate
\usepackage{multicol}
\usepackage[many]{tcolorbox}
\usepackage{hyperref}
\usepackage{arydshln}

%\renewcommand{\ttdefault}{lmtt} %Schriftart wird geändert zu Latin Modern Typewriter


\hypersetup{hidelinks,
			colorlinks=true,
			linkcolor=black,
			citecolor=miudiutex,
			urlcolor=miugray2
			}

\setlength{\parindent}{0pt}
\setcounter{section}{0}


\titlespacing*{\section}{0pt}{20mm}{6mm}
\titlespacing*{\subsection}{0pt}{6mm}{3mm}
\titlespacing*{\subsubsection}{0pt}{3mm}{0mm}



%\definecolor{miudiutex}{HTML}{2f3d39}
\definecolor{miudiutex}{HTML}{1d2926}
\definecolor{miugray}{HTML}{b4cccf}
\definecolor{miugray2}{HTML}{4d7365}
\definecolor{red}{HTML}{cf1f5c}
\definecolor{green}{HTML}{00A36C}
\definecolor{delphinidin}{HTML}{315794}
\definecolor{pelargonidin}{HTML}{b33807}
\definecolor{cyanidin}{HTML}{ba0746}
\definecolor{petunidin}{HTML}{b56fed}



\raggedcolumns %wichtig für Multicolumns


%================================
% SYMBOLS
%================================

\newcommand{\RA}{$\rightarrow$~}
\newcommand{\RL}{$\rightleftharpoons$~}
\newcommand{\HARP}{\ding{226}~}
%$\xrightarrow{2}$ %Pfeil mit 2 drüber


%================================
% SHORT COMMANDS (BOXES ; ETC)
%================================


\newcommand{\notiz}[1]{\begin{notizz}{}{}\begin{center}
			{\color{miugray2!50!miudiutex}#1} \end{center}\end{notizz}}
\newcommand{\zit}[2]{\begin{zitat*}{#1}{}\small #2 \end{zitat*}}

\newcommand{\frage}[2]{\begin{qst*}{#1}{}\small #2 \end{qst*}}

\newcommand{\unten}{\end{center} \tcblower \begin{center}}


\newcommand*\circled[1]{\tikz[baseline=(char.base)]{
            \node[shape=circle,draw,inner sep=0.5pt,miudiutex] (char) {\tiny #1};}}
            
\newcommand{\miusquare}{{\color{miugray2!50}\scriptsize $\blacksquare$}~}
\newcommand{\miusquarp}{{\color{petunidin!50}\scriptsize $\blacktriangleright$}~}
\newcommand{\niu}{\item[\miusquare]}
\newcommand{\nip}{\item[\miusquarp]}

%%%%% ABSAETZE

\newcommand{\pps}{\par \vspace*{1mm}}
\newcommand{\pp}{\par \vspace*{3mm}}
\newcommand{\ppbig}{\par \vspace*{8mm}}
\newcommand{\pphuge}{\par \vspace*{10mm}}

%================================
% FRAGE BOX
%================================

\tcbuselibrary{theorems,skins,hooks}
\newtcbtheorem{qst}{Frage:}
{%
	enhanced,
	breakable,
	colback = gray!15!white,
	frame hidden,
	boxrule = 0sp,
	borderline west = {2.5pt}{0pt}{red!70},
	sharp corners,
	detach title,
	before upper = \tcbtitle\par\smallskip,
	coltitle = black,
	fonttitle = \bfseries\sffamily,
	description font = \mdseries,
	separator sign none,
	segmentation style={solid, gray!50!black},
}
{th}




%================================
% ZITAT BOX
%================================

\tcbuselibrary{theorems,skins,hooks}
\newtcbtheorem{zitat}{Zitat}
{%
	enhanced,
	breakable,
	colback = gray!15!white,
	frame hidden,
	boxrule = 0sp,
	borderline west = {2.5pt}{0pt}{black!70},
	sharp corners,
	detach title,
%	before upper = \tcbtitle\par\smallskip,
	coltitle = gray!50!black,
	fonttitle = \bfseries\sffamily,
	description font = \mdseries,
%	separator sign none,
	segmentation style={solid, gray!50!black},
}
{th}



%================================
% SIMPLE SCHWARZER RAHMEN BOX
%================================


\newcommand{\boxfr}[1]{
	\begin{tcolorbox}[
	drop shadow southeast,
	enhanced,
	colframe=black!50,
	colback=white,
	boxrule = 1pt, 
%	toprule=0pt, bottomrule=0pt,rightrule=0pt,leftrule=0pt,
	boxsep=5pt,
	arc=1pt,
	]
	\begin{center}
	#1
	\end{center}
	\end{tcolorbox}
}

\definecolor{miudiutex}{HTML}{1d2926}
\definecolor{miugray}{HTML}{b4cccf}
\definecolor{miugray2}{HTML}{4d7365}
\definecolor{white2}{HTML}{f5f7f8}
\definecolor{red}{HTML}{cf1f5c}
\definecolor{green}{HTML}{00A36C}
\definecolor{delphinidin}{HTML}{315794}
\definecolor{uniblau}{HTML}{004e9f}
\definecolor{unigelb}{HTML}{fcba00}
\definecolor{unigrau}{HTML}{909085}
\definecolor{pelargonidin}{HTML}{b33807}
\definecolor{cyanidin}{HTML}{ba0746}
\definecolor{paeonidin}{HTML}{d15ca6}
\definecolor{petunidin}{HTML}{b56fed}
\definecolor{malvidin}{HTML}{8a2d8c}

\newcommand{\metermilli}{\mskip3mu \text{mm}}
\newcommand{\cm}{\mskip3mu \text{cm}}
\newcommand{\m}{ \text{m}}
\newcommand{\lite}{ \text{l}}
\newcommand{\kg}{ \text{kg}}
\newcommand{\g}{ \text{g}}
\newcommand{\km}{ \text{km}}
\newcommand{\h}{ \text{h}}
\newcommand{\s}{ \text{s}}
\newcommand{\tonne}{ \text{t}}
\usepackage[utf8]{inputenc}
\usepackage[T1]{fontenc}
\usepackage{lmodern}

\usepackage[
    activate={true,nocompatibility},
    final,
    tracking=true,
    kerning=true,
    spacing=true,
    factor=1100,
    stretch=30,
    shrink=20
]{microtype}
\usepackage{pdfpages}
\usepackage{awesomebox}
\usepackage{marginnote}
\tcbuselibrary{documentation}
\usepackage{sidenotes}
\usepackage{import}
\usepackage{xifthen}
\usepackage{transparent}
\usepackage[table]{xcolor}


\newcommand{\abbicon}{%
  \protect\tikz[baseline=0.2mm,scale=0.8]{%
    % Das blaue Quadrat mit abgerundeten Ecken
    \fill[uniblau, rounded corners=1.2pt] (0,0) rectangle (0.8em, 0.8em);
    % Der weiße Kreis in der Mitte
    \fill[white] (0.4em, 0.4em) circle (0.12em);
  }%
} 
\usepackage[labelfont={bf},font={color=miudiutex,small},figurename={\abbicon $~$Abbildung}]{caption}


\newcommand{\bckgrnd}{
\AddToHookNext{shipout/background}{
\begin{tikzpicture}[remember picture, overlay]
 \draw[draw=none,fill=uniblau!70, shift={(current page.north west)}] (-3,0) rectangle (23,-3.75);
 \draw[draw=none,fill=miugray, shift={(current page.north east)}] (-7,0) rectangle (3,-3.75);
\end{tikzpicture}
}
}



\newcommand{\incfig}[2]{%
\def\svgwidth{#2\columnwidth}
\import{C://LaTeX-Files/pngs/}{#1.pdf_tex}
}

\renewcommand{\textbf}[1]{\textsf{{\bfseries {#1}}}}



\titleformat{\section}[hang]
		{\LARGE \color{uniblau!75!black} \sffamily \bfseries}
		{\thesection}
		{6mm}
		{}
		[]
\titleformat{\subsection}[hang]
		{\large \sffamily \bfseries \color{black}}
		{\thesubsection}
		{4mm}
		{{\color{miugray}$\blacksquare ~$}}

\titleformat{\subsubsection}[block]
		{\normalsize \bfseries \sffamily \color{miudiutex}}%
		{\normalsize \color{black} \thesubsubsection}
		{2mm}
		{{\color{miugray}$\blacksquare ~$}}
		[ \color{black}]
		
\titleformat{\paragraph}
		{\normalsize \bfseries \sffamily \color{black}}
		{\footnotesize \theparagraph}
		{1mm}
		{}

\pagestyle{empty}
\titlespacing*{\section}{0pt}{20mm}{12mm}
\titlespacing{\section}{0pt}{20mm}{12mm}

\titlespacing*{\subsection}{0pt}{14mm}{4mm}
\titlespacing{\subsection}{0pt}{14mm}{4mm}

\titlespacing*{\subsubsection}{0pt}{6mm}{3mm}
\titlespacing{\subsubsection}{0pt}{6mm}{3mm}

\titlespacing*{\paragraph}{0pt}{6mm}{2mm}
\titlespacing{\paragraph}{0pt}{6mm}{2mm}


\let\oldmarginfigure\marginfigure
\let\endoldmarginfigure\endmarginfigure

\let\oldmarginfigure\marginfigure
\let\endoldmarginfigure\endmarginfigure

\renewenvironment{marginfigure}[1][0pt] % [1] steht für 1 Argument, [0pt] ist der Standardwert
  {\oldmarginfigure[#1]\captionsetup{labelfont={sf,bf,color=uniblau!75!black},font={color=miudiutex,small}}}
  {\endoldmarginfigure}
  
  
  
  
\renewcommand{\labelitemi}{\color{uniblau!70}$\bullet$}
\renewcommand{\labelitemii}{\color{uniblau!70}$\cdot$}
\newcommand{\plmtsquare}{{\color{uniblau!70}$\blacksquare~$}}
\newcommand{\splmt}[1]{{\color{uniblau!75!black} \sffamily #1}}
\newcommand{\fplmt}[1]{{\color{uniblau!75!black} \sffamily \bfseries #1}}

\newcommand{\antwort}[1]{
\begin{tcolorbox}[
	enhanced,
	enforce breakable,
	standard jigsaw,
	boxrule=0.7pt,
	colframe=black,
	sharp corners,
	boxsep=5pt,
	title=\bfseries \sffamily Antwort,
	opacityback=0,
	coltitle=miudiutex,
	attach boxed title to top text left={yshift=-\tcboxedtitleheight/2},
	boxed title style={colback=white,colframe=white},
	bottomrule at break=0pt,
	toprule at break=0pt,
	pad at break=16mm,
]
	\begin{center}
	#1
	\end{center}
	\end{tcolorbox}
}




\rowcolors{2}{miugray!50}{miugray!50}
\arrayrulecolor{white} % Weiße Gitterlinien
\setlength{\arrayrulewidth}{1.5pt} % Etwas dickere Linien
\renewcommand{\arraystretch}{1.4} % Mehr Zeilenhöhe für bessere Lesbarkeit}


%\setlist[itemize]{itemsep=0mm}
\setlist[enumerate]{itemsep=0mm}
\usepackage[table]{xcolor}
\usepackage{pgfplots}
    \pgfplotsset{
        every axis post/.style={
	yticklabel=\empty,
	xticklabel=\empty,
	y label style={at={(axis description cs:-0.2,1.1)}},
	x label style={at={(axis description cs:1.05,0.2)}},
	ticks=none,
	width=6cm,
	axis lines=middle,
%	xlabel near ticks,
        },
    }
\usetikzlibrary{positioning, arrows.meta, calc, patterns}

\usepackage{relsize}

\newif\ifshowme
\showmetrue




\begin{document}
\newgeometry{top=1.8cm,bottom=4cm,left=1.5cm,right=9.5cm,marginparsep=-4.00cm,marginparwidth=6.75cm}


\bckgrnd

\begin{center}
\LARGE
\color{white}
\textbf{Tafelanschrift} \par \vspace*{2mm}
\large \color{miugray} \textbf{Dienstag, 21.04.2026}
\end{center}
\par 
\vspace*{6mm}


\color{miudiutex}

\section*{Laminare Rohrströmung mit Reibung}
\small
\begin{marginfigure}
\resizebox{1.2\textwidth}{!}{
\begin{tikzpicture}[font=\sffamily]

    % Überschrift und Text
%    \node[anchor=west] at (-4, 4.5) {\Large \underline{Laminare Rohrströmung mit Reibung (für lange Rohre)}};
%    \node[anchor=west] at (-3, 3.5) {Annahme: inkompressibles, newtonsches Fluid};

    % Linke Textspalte
%    \node[anchor=west, text width=5cm] at (-5, 1.5) {
%        $\bullet$ Antrieb für eine\\ 
%        \hspace*{0.4cm} Rohrströmung: $\dfrac{dp}{dx}$\\[2mm]
%        \hspace*{0.4cm} Druckgradient\\[5mm]
%        $\bullet$ Reibungswiderstand\\
%        \hspace*{0.4cm} kommt durch Viskosität\\
%        \hspace*{0.4cm} des Fluids zustande
%    };

    % Das Rohr
    \draw[thick] (0, 1.5) -- (10, 1.5);
    \draw[thick] (0, -1.5) -- (10, -1.5);
    
    % Linker Rohrquerschnitt (Ellipse)
    \draw[thick] (0,0) circle [x radius=0.6, y radius=1.5];
    \draw[->, >=Stealth] (0,0) -- (0.5, 1.4) node[midway, left] {$R$};

    % Zylindrisches Fluidelement
    \begin{scope}[shift={(4,0)}]
        \draw[thick] (0,0) circle [x radius=0.3, y radius=0.75];
        \draw (0,0) -- (0.2, 0.7) node[midway, left, scale=0.8] {$r$};
        \node at (0,-0.3) {$A$};
        \draw[thick] (0, 0.75) -- (2, 0.75);
        \draw[thick] (0, -0.75) -- (2, -0.75);
        \draw[thick, dashed] (2,0) circle [x radius=0.3, y radius=0.75];
        
        % Kräfte am Element
        \draw[->, thick, >=Stealth] (-1, 0) -- (-0.1, 0) node[pos=0, above] {$F_{p_1}$};
        \draw[<-, thick, >=Stealth] (2.1, 0) -- (3, 0) node[pos=1, above] {$F_{p_2}$};
        
        % Längenmaß dx
        \draw[|<->|] (0, -1) -- (2, -1) node[midway, below] {$dx$};
    \end{scope}

    % Geschwindigkeitsprofil (Parabel)
    \begin{scope}[shift={(8,0)}]
        \draw[dashed] (0, -1.5) -- (0, 1.5);
        \draw[domain=-1.5:1.5, smooth, variable=\y, thick, uniblau] plot ({1.5*(1-(\y/1.5)^2)}, \y);
        
        \foreach \y in {-1.2,-0.9,-0.6,-0.3,0,0.3,0.6,0.9,1.2}
            \draw[->, >=Stealth, uniblau!50!miugray] (0,\y) -- ({1.5*(1-(\y/1.5)^2)}, \y);
            
        \node[anchor=south west] at (0, 1.5) {$v(r)$};
        \node[anchor=west, scale=0.8] at (1.5, 0) {$v_{max}$};
    \end{scope}

    % Legende unten
    \node[anchor=west] at (3.5, -2.5) {$r$: Radius des zylindrischen Fluidelements};
    \node[anchor=west] at (3.5, -3.1) {$A$: Stirnfläche \hspace{3.8cm} $(A = \pi r^2)$};

\end{tikzpicture}
}
\caption{Schematische Darstellung; Annahme: inkompressibles, newtonsches Fluid}
\end{marginfigure}


\subsubsection*{Druckkraft}
Der Antrieb für eine Rohrströmung ergibt sich aus dem Druckgradienten $\frac{dp}{dx}$
\begin{align*}
F_{p_1} &= p \cdot \pi r^2 \\
F_{p_2} &= (p + dp) \cdot \pi r^2
\end{align*}

Effektiv wirkende Kraft $F_p$
\begin{align*}
F_p = F_{p_1} - F_{p_2} = -\pi r^2 \, dp
\end{align*}


\subsubsection*{Reibungswiderstand}
Der Reibungswiderstand kommt durch die Viskosität des Fluids zustande.
\begin{align*}
\tau &= \eta \cdot \frac{dv(r)}{dr} \\
\tau &= \frac{F}{A} \\ \\
F_R &= \tau \cdot A \\
&= \tau \cdot 2\pi r \cdot dx \\
&= \eta \cdot \frac{dv(r)}{dr} \cdot 2\pi r \cdot dx
\end{align*}

\subsubsection*{Gleichsetzen}
Gleichsetzen der Druckkraft und der Reibungskraft (stationäre Strömung)
\marginnote[]{\small $\mathlarger{\frac{dv(r)}{dr}}$ = Geschwindigkeitsgradient der Strömung in radialer Richtung \pp
$\mathlarger{\frac{dp}{dx}}$ = Druckgradient in axialer Richtung}
\begin{align*}
F_R &= F_p \\
\eta \cdot \frac{dv(r)}{dr} \cdot 2\pi r \cdot dx &= -\pi r^2 \cdot dp \\
\frac{dv(r)}{dr} &= -\frac{1}{2\eta} \cdot \frac{dp}{dx}\\
\end{align*}


\subsubsection*{Integration}\vspace*{-6mm}

\marginnote[]{\small \color{uniblau}
Strömungsgeschwindigkeit für r = R: 0
\begin{align*}
-\frac{1}{4\eta} \cdot \frac{dp}{dx} \cdot R^2 + C = 0 \\
C = \frac{1}{4\eta} \cdot \frac{dp}{dx} \cdot R^2
\end{align*}
}[1.5cm]
\begin{align*}
v(r) &= \int -\frac{1}{2\eta} \cdot \frac{dp}{dx} \cdot r \cdot dr \\
v(r) &= -\frac{1}{4\eta} \cdot \frac{dp}{dx} \cdot r^2 \color{uniblau}{+~ C} \\
&= -\frac{1}{4\eta} \cdot \frac{dp}{dx} \cdot r^2 + \frac{1}{4\eta} \cdot \frac{dp}{dx} \cdot R^2
\end{align*}

Gesetz von Stokes:
\begin{align*}
v(r) &= -\frac{1}{4\eta} \cdot \frac{dp}{dx} \cdot (R^2 - r^2) \\
v(r) &= \underbrace{-\frac{R^2}{4\eta} \cdot \frac{dp}{dx}}_{=v_{max}} \cdot \left[ 1 - \left(\frac{r}{R}\right)^2 \right]
\end{align*}
\subsubsection*{Maximale Geschwindigkeit}\vspace*{-6mm}
\begin{align*}
v_{max} &= v(r=0) \\
&= - \frac{R^2}{4\eta} \cdot \frac{dp}{dx} \cdot \left[ 1 - \left(\frac{0}{R}\right)^2 \right] \\
&= - \frac{R^2}{4\eta} \cdot \frac{dp}{dx}
\end{align*}






\end{document}
